\scanpage{261}
\mspage{2}

intervalles standard de longueur donnée $n$ correspondent biunivoquement aux syst. de signes \deleted{\unsure{$\pm$}} $+$ \inserted{ou $-$} sur les $n$ intervalles de transition $[i-1,i]$ ($1\leq i\leq n$), il y en a donc $2^n$. Le sommet $0$ est l'\uline{origine} \inserted{ou la \uline{source}} de l'intervalle, le sommet $n$ l'\uline{extrémité} ou le \uline{but} de l'intervalle.

Un intervalle standard est \uline{orienté}, du fait qu'on a la relation d'ordre \deleted{\uline{totale}} naturelle sur \deleted{\unsure{l}}l'ens. de ses sommets, \inserted{savoir celle des entiers de $[0,n]$.} Quand on parlera de la relation d'ordre sur $I$, qu'un sommet $i$ en ``\uline{précède}'' ou ``\uline{suit}'' un autre $j$ ($\unsure{i|v}<j$ ou $i\unsure{>|{\geq}}j$ respectivement), c'est de cette relation d'ordre qu'on voudra parler. Celle associée à la structure de catégorie ne sera \unsure{jamais} regardée comme une relation d'ordre.

Il y a également une relation d'ordre totale naturelle sur l'ens. des flèches de transition $\underline{u}_i$ \inserted{$(1\leq i\leq n)$}, donnée par \deleted{\unsure{la}} la relation d'ordre des entiers $i$, $1\leq i\leq n$, qui les indexent. On dira qu'une flèche $\unsure{i|\prime}$ de tr. \uline{précède} ou \uline{suit} une autre, ou \deleted{\unsure{la}} qu'elle est \uline{successeur} \inserted{\uline{immédiat}} ou \uline{prédécesseur immédiat} d'une autre (si $i=j+1$, ou $i=j-1$).

Si on a un ensemble d'entiers

\[E\subset\mathbb{Z},\]

on \deleted{\unsure{dira}} considérera dans $E$ la relation d'équivalence engendrée par la relation entre deux éléments $(i,j)$ de $E$ : $j$ est successeur immédiat de $i$ \uline{dans} $\deleted{\unsure{j}}\mathbb{Z}$ (i.e. $j=i+1$). Deux éléments \inserted{distincts $i,j$} de $E$ \deleted{qui} sont congrus pour cette relation \unsure{sss|ssi} \deleted{\unsure{on a}} l'intervalle d'entiers qu'ils déterminent est $\subset E$, donc

\[\text{si }i<j,\quad i\text{ et }j\text{ sont ``\uline{connectés}'' dans }E\text{ ssi }[i,j]\subset E.\]

% archive sheet number: 260
